\section{From a full configuration to Gaussian-smoothed marginals}\label{sec:gaussian}

We first fix a finite labelled multiset $\mathbf z=(z_1,\ldots,z_N)$.
Labels, rather than geometric positions, distinguish multiplicity copies.
All densities in this section use planar Lebesgue measure $\dd A$ in
each coordinate. The construction is an observation law for a given
configuration, not a hypothesis that canonical zeros are independent.
The factorial-measure convention agrees with the standard correlation
framework \cite[Section 2]{Borodin2009DPP}, extended here to labelled
multiple zeros before smoothing.

\subsection{The full joint density and its marginals}

For $\epsilon>0$ define the normalized planar Gaussian
\begin{equation}\label{eq:gaussian}
 G_\epsilon(u)=\frac1{2\pi\epsilon^2}
     \exp\!\left(-\frac{|u|^2}{2\epsilon^2}\right),
 \qquad \int_\C G_\epsilon\dd A=1.
\end{equation}
Randomly permute the labels and perturb each observed position by an
independent Gaussian with this density. The full joint probability
density is
\begin{equation}\label{eq:full-joint}
 p_{N,\epsilon}(u_1,\ldots,u_N)
 =\frac1{N!}\sum_{\sigma\in\mathfrak S_N}
      \prod_{a=1}^N G_\epsilon(u_a-z_{\sigma(a)}).
\end{equation}
Write $(N)_k=N!/(N-k)!$ for $1\le k\le N$.
Integrating the last $N-k$ variables gives
\begin{equation}\label{eq:joint-marginal}
 p^{(k)}_{N,\epsilon}(\mathbf u)
 =\frac1{(N)_k}\sum_{j_1,\ldots,j_k\ \mathrm{distinct}}
            \prod_{a=1}^k G_\epsilon(u_a-z_{j_a}).
\end{equation}
Indeed the integrated factors each have mass one and each ordered
$k$-tuple occurs $(N-k)!$ times. The associated factorial correlation
function is $R_{k,\epsilon}=(N)_k p^{(k)}_{N,\epsilon}$. Consequently
\begin{equation}\label{eq:fixed-N-consistency}
 \int_{\C^k}R_{k,\epsilon}\dd A^{\otimes k}=(N)_k,
 \qquad
 \int_\C R_{k,\epsilon}(\mathbf u)\dd A(u_k)
       =(N-k+1)R_{k-1,\epsilon}(u_1,\ldots,u_{k-1}).
\end{equation}
These are genuine marginals of a full joint law. They are also
Gaussian convolutions of the deterministic factorial zero measures.
Neither interpretation assigns a new probability distribution to the
unsmoothed canonical zeros.

\subsection{Smooth spatial localization and weighted normalization}

Choose a compactly supported $C^2$ function $0\le W\le1$ on $\C$.
For the zeros of a nonzero entire function only finitely many weights
$w_j=W(z_j)$ are nonzero. Our primary localized objects are
\begin{align}
 \alpha_k^W&=\sum_{j_1,\ldots,j_k\ \mathrm{distinct}}
       \left(\prod_{a=1}^k w_{j_a}\right)
       \delta_{(z_{j_1},\ldots,z_{j_k})},\label{eq:weighted-factorial}\\
 R_{k,\epsilon}^W(\mathbf u)
   &=\int\prod_{a=1}^kG_\epsilon(u_a-z_a)\dd\alpha_k^W(\mathbf z)
     =\sum_{j_1,\ldots,j_k\ \mathrm{distinct}}
          \prod_{a=1}^k w_{j_a}G_\epsilon(u_a-z_{j_a}).\label{eq:full-correlation}
\end{align}
Set $R_{0,\epsilon}^W=1$ and $M_W=\sum_jw_j$.
The window acts on the true zero locations, before the observation
noise. Thus $R_{k,\epsilon}^W$ has Gaussian tails in observed coordinates
even though its source test is compactly supported.

There is a consistent point-process interpretation: retain label $j$
independently with probability $w_j$, and then perturb each retained
location as above. If $M$ is the random retained count, conditioning on
$M=m$ and uniformly labelling the retained observations gives symmetric
probability marginals $p_m^{(k)}$. Their mixture satisfies
\begin{equation}\label{eq:random-N-consistency}
 R_{k,\epsilon}^W
    =\sum_{m\ge k}(m)_k\Pr(M=m)p_m^{(k)},\qquad
 Z_k:=\int R_{k,\epsilon}^W=k!e_k(w_1,\ldots,w_N)\le M_W^k.
\end{equation}
The equality follows by conditioning each ordered retained tuple;
the inequality restores the tuples with repeated labels. If $Z_k>0$,
$R_{k,\epsilon}^W/Z_k$ is a probability density for a weighted ordered
tuple. These separately normalized densities need not be mutually
consistent marginals of a fixed-size law: integrating the last coordinate
in \eqref{eq:full-correlation} leaves the factor
$M_W-\sum_{a<k}w_{j_a}$ for each remaining tuple, not a constant.
The fixed-$N$ rule \eqref{eq:fixed-N-consistency} applies only to the
unweighted retained set, or conditionally to a declared fixed count.

This distinction also separates localization from reweighting an
ensemble joint density by $\prod_jW(z_j)$ and renormalizing. That latter
operation changes the configuration law and its partition function.
Here the canonical configuration is fixed; a Vandermonde ensemble would
be a separate model, not a consequence of the spatial window.

\begin{proposition}[The weighted pair and its self-label correction]\label{prop:weighted-pair}
Define
\begin{equation}\label{eq:r-and-d}
 r_\epsilon^W(u)=\sum_j w_jG_\epsilon(u-z_j),\qquad
 d_\epsilon^W(u,v)=\sum_jw_j^2G_\epsilon(u-z_j)G_\epsilon(v-z_j).
\end{equation}
Then the full two-zero correlation function is
\begin{equation}\label{eq:pair-rd}
 R_{2,\epsilon}^W(u,v)
       =r_\epsilon^W(u)r_\epsilon^W(v)-d_\epsilon^W(u,v),
 \qquad Z_2=M_W^2-\sum_jw_j^2.
\end{equation}
At a location with $m$ multiplicity copies, its coincident-label-group
contribution is $m(m-1)W(z)^2G_\epsilon(u-z)G_\epsilon(v-z)$.
\end{proposition}
\begin{proof}
Expand the product of the two sums defining $r$. Its $j=k$ terms
are exactly $d$, with weight $w_j^2$. The remaining terms are ordered
distinct labels, including distinct labels at the same position.
Integrating both Gaussian factors proves the mass formula.
\end{proof}

The square on $w_j$ in $d$ is essential. The ordinary second moment
of the randomly thinned counting measure instead has a diagonal
expectation weighted by $w_j$, because a retained label is counted once.
Equation~\eqref{eq:pair-rd} subtracts the repeated index in a product
of first intensities, which has weight $w_j^2$.

\subsection{A generating functional for all orders}

For a bounded complex test $f$, let
\begin{equation}\label{eq:full-pgfl}
 \mathcal G_{W,\epsilon}[f]
  =\prod_j\left(1+w_j\int_\C f(u)G_\epsilon(u-z_j)\dd A(u)\right).
\end{equation}
This is the expectation of $\prod_{u\ \mathrm{observed}}(1+f(u))$ under
the thinning-and-smoothing law. Expanding the finite product yields
\begin{equation}\label{eq:pgfl-expansion}
 \mathcal G_{W,\epsilon}[f]
  =\sum_{k=0}^N\frac1{k!}
      \int_{\C^k}\prod_{a=1}^k f(u_a)
                       R_{k,\epsilon}^W(\mathbf u)\dd A^{\otimes k}.
\end{equation}
If $\|f\|_\infty\le q$, the absolute tail after order $K$ is at most
\begin{equation}\label{eq:order-tail-full}
 \sum_{k>K}\frac{(qM_W)^k}{k!}
       \le \ee^{qM_W}\frac{(qM_W)^{K+1}}{(K+1)!}.
\end{equation}
For the last inequality put $k=K+1+j$ and use
$(K+1+j)!\ge(K+1)!j!$. This bound pays the order truncation
without identifying different orders or postulating a closed hierarchy.

\subsection{Smoothing flow, bias, and localization tails}

Put $\tau=\epsilon^2/2$. Direct differentiation of
$G_{\sqrt{2\tau}}=(4\pi\tau)^{-1}\ee^{-|u|^2/(4\tau)}$ gives
$\partial_\tau G_{\sqrt{2\tau}}=\Delta G_{\sqrt{2\tau}}$.
The finite sums therefore satisfy
\begin{equation}\label{eq:smoothing-flow}
 \partial_\tau R_{k,\sqrt{2\tau}}^W
    =\sum_{a=1}^k\Delta_{u_a}R_{k,\sqrt{2\tau}}^W,
 \qquad R_{k,\sqrt{2\tau}}^W\dd A^{\otimes k}
      \rightharpoonup\alpha_k^W\quad(\tau\downarrow0).
\end{equation}
Weak convergence follows by testing each Gaussian approximate identity.
This is an auxiliary observation-smoothing time. It is distinct from
Newman time $\lambda$ in \eqref{eq:heat}.

For $F\in C_b^2(\C^k)$, with $\mathcal L_k=\sum_a\Delta_{u_a}$,
the Gaussian semigroup and its sup-norm contraction give
\begin{equation}\label{eq:smoothing-bias}
 \left|\int F R_{k,\epsilon}^W\dd A^{\otimes k}
                       -\alpha_k^W(F)\right|
 \le\frac{\epsilon^2}{2}\,Z_k\|\mathcal L_kF\|_\infty.
\end{equation}
Indeed $P_\tau F-F=\int_0^\tau P_t\mathcal L_kF\dd t$; integrate
against the positive measure $\alpha_k^W$. This is a weak bias estimate,
not uniform convergence of densities to an atomic measure.

For a countable weighted configuration with $M_W<\infty$, all the
preceding nonnegative sums converge. If a truncation omits total weight
$T$, let $R_{k,\epsilon}^{W,\mathrm{tr}}$ denote the retained sum and put
$A_\epsilon=(2\pi\epsilon^2)^{-1}$. Then
\begin{equation}\label{eq:spatial-tail}
 \begin{split}
 0&\le R_{k,\epsilon}^W-R_{k,\epsilon}^{W,\mathrm{tr}}
          \le k A_\epsilon^kM_W^{k-1}T,\\
 \|R_{k,\epsilon}^W-R_{k,\epsilon}^{W,\mathrm{tr}}\|_{L^1}
          &\le kM_W^{k-1}T.
 \end{split}
\end{equation}
Every omitted tuple has an omitted label in at least one of $k$ slots;
sum over the choice of slot and restore repeated labels to bound the
remaining $k-1$ sums by $M_W^{k-1}$. This also proves the $L^1$ claim
after integrating the Gaussians. Section~\ref{sec:gaussian-window}
pays the corresponding canonical mass and tail estimates for an
untruncated Gaussian physical window. The source theorem in
Section~\ref{sec:full-pair} uses a compact physical window and Gaussian
observation smoothing.
